\section{Method}

\subsection{System Stack and Problem Setting}

We consider suspended-payload UAV transport under wind, where the vehicle must
move a cable-suspended load to a target while limiting tracking degradation,
payload swing, and arrival failures. The suspended payload couples the UAV
motion to load dynamics, so execution aggressiveness can change the effective
difficulty of the same nominal trajectory. Under strong wind, a command that
is feasible in one run can become too aggressive or poorly timed in another.

The underlying system follows an AutoTrans-like stack composed of a
payload-aware planner, a payload MPC, and an SO(3) controller
\cite{Li2023AutoTrans,UrbinaBrito2021PredictivePayloadTransport,Lee2010GeometricTrackingSE3}.
The planner generates nominal motion references, the payload MPC tracks these
references while accounting for the load, and the SO(3) controller provides the
low-level vehicle control interface. In this paper, these core stack
components are kept unchanged.

\subsection{Execution-Governor Interface}

The execution governor publishes two scalar adaptation signals:
speed scale and acceleration scale. These signals are
consumed by the command-adaptation interface and modify the aggressiveness of
the executed motion command. A scale near 1.0 preserves the nominal command
more closely, whereas smaller scales reduce the commanded speed or acceleration
and make the executed motion more conservative.

This interface is intentionally stack-compatible. The planner continues to
produce the nominal reference, the payload MPC continues to solve the tracking
problem, and the SO(3) controller remains the low-level controller. The governor
only changes the execution scale exposed through the command-adaptation
interface. This design connects to command and reference governance concepts
\cite{Garone2017ReferenceCommandGovernorsSurvey,Garone2016ExplicitReferenceGovernor,Li2021ActionGovernor},
but the implementation here is an empirical execution policy rather than a
certified reference governor.

\begin{figure}[t]
  \centering
  \resizebox{\linewidth}{!}{\begin{tikzpicture}[
  font=\sffamily\scriptsize,
  >=Latex,
  block/.style={
    draw=black!72,
    fill=black!3,
    rounded corners=1pt,
    minimum width=1.55cm,
    minimum height=0.62cm,
    align=center,
    inner sep=2.4pt
  },
  wideblock/.style={
    block,
    minimum width=1.95cm
  },
  flow/.style={-Latex, line width=0.58pt, draw=black!82},
  adapt/.style={-Latex, line width=0.62pt, draw=black!78, dashed},
  sense/.style={-Latex, line width=0.54pt, draw=black!66, densely dotted},
  group/.style={draw=black!48, rounded corners=2pt, dashed, inner sep=4.8pt}
]

\node[block] (planner) at (0.00,1.35) {Planner /\\reference\\generator};
\node[block] (mpc) at (2.05,1.35) {Payload\\MPC};
\node[block] (so3) at (3.95,1.35) {SO(3)\\controller};
\node[wideblock] (plant) at (6.10,1.35) {UAV +\\suspended\\payload};

\node[
  group,
  fit=(planner)(mpc)(so3),
  label={[font=\sffamily\tiny, yshift=1pt]above:unchanged inner stack}
] (stack) {};

\node[wideblock] (risk) at (0.85,-0.58) {Risk input /\\warning scores};
\node[wideblock, minimum width=2.10cm] (governor) at (3.25,-0.58)
  {Risk-conditioned\\execution governor\\{\tiny empirical governor}};
\node[wideblock, minimum width=2.05cm] (interface) at (5.70,-0.58)
  {Command-adaptation\\interface\\{\tiny speed scale}\\{\tiny acceleration scale}};
\node[wideblock, minimum width=1.78cm] (logs) at (6.10,-1.78)
  {Diagnostics /\\logs};

\draw[flow] (planner) -- (mpc);
\draw[flow] (mpc) -- (so3);
\draw[flow] (so3) -- (plant);

\draw[sense] (plant.south) -- ++(0,-0.38) -| (logs.north);
\draw[sense] (plant.south west) .. controls (4.70,0.38) and (2.25,0.12) ..
  (risk.north east);
\draw[flow] (risk) -- (governor);
\draw[flow] (governor) -- (interface);
\draw[adapt] (interface.north) .. controls (5.65,0.14) and (5.00,0.60) ..
  ($(so3.south)!0.58!(plant.south)$);
\node[font=\sffamily\tiny, align=center, fill=white, inner sep=1.2pt]
  at (5.18,0.25) {execution\\command path};
\draw[sense] (logs.west) -| (risk.south);

\end{tikzpicture}
}
  \caption{Stack-compatible risk-conditioned execution governance. The
  planner, payload MPC, and SO(3) controller remain unchanged, while an external
  governor adapts execution aggressiveness through speed scale and
  acceleration scale. The governor is empirical and is not a
  certified safety filter.}
  \label{fig:architecture}
\end{figure}

\subsection{Risk-Conditioned Governor}

The learned governor uses risk-conditioned decision logic to choose execution
scales. At a high level, the governor receives short-horizon and long-horizon
risk estimates, denoted here as $r_3$ and $r_5$, together with execution
context such as the active method mode and command-adaptation state. The
short-horizon risk captures imminent instability or tracking stress, while the
long-horizon risk captures a broader upcoming window where conservative
execution may be useful.

In the current implementation, these scores are exposed as
3 s and 5 s risk scores. The adapter subscribes
to odometry, payload odometry, wind annotation, goal, and trajectory-context
streams, builds early episode features, and evaluates exported
JSON-serialized logistic regression models trained with
the strict-invalid label as the positive class. The audited local JSON
metadata records a 90-row training dataset with 51 strict-valid and 39
strict-invalid rows. The 3 s model uses 16 wind, target, and
3 s early-window features, while the 5 s model uses 27 features spanning the
same base fields plus 3 s and 5 s early-window dynamics; the
full feature schema is kept for supplementary material. Unavailable scores are
published as a sentinel value of -1 and do not trigger risk reductions. These
values are used as empirical strict-invalid warning scores, not as calibrated
physical probabilities; the current implementation does not include confidence
or OOD rejection, and inference latency has not been measured separately from the
configured 5 Hz publish rate.

The governor maps these signals to a pair of scale values. In low-risk
conditions, it keeps execution less conservative so that the system can
continue progressing toward the goal. When risk increases, it reduces the speed
and acceleration scales to avoid overly aggressive execution. When the
long-horizon risk is high enough, the governor may apply a stronger reduction.
The scale output is rate-limited so that the command does not switch abruptly
between aggressive and conservative modes.

\subsection{Current Protagonist: Risk Adapter v1}

The current paper protagonist is Risk Adapter v1, the balanced learned /
risk-conditioned execution governor. Its policy settings are:

\begin{itemize}
  \item 3 s risk threshold: 0.5
  \item 5 s risk threshold: 0.5
  \item 5 s hard threshold: 0.7
  \item 3 s soft scale: 0.75
  \item 5 s soft scale: 0.65
  \item 5 s hard scale: 0.60
  \item scale rate limit: 0.5 per second
\end{itemize}

Risk Adapter v1 is selected as the tentative balanced protagonist because it
gives the strongest observed cross-protocol balance in the current evaluation.
It achieves 25/30 strict-valid runs in the single-goal mission protocol and
23/30 in the goal-reissue stress protocol, yielding the highest mean valid
count (24.0/30), highest worst-protocol valid count (23/30), and lowest total
regret (2) relative to the observed protocol oracles.

The later Risk Adapter v2.1 variant remains important, but it should not be
presented as the final method. It ties Risk Adapter v1 in the
single-goal mission protocol at 25/30, but it drops to 20/30 under goal-reissue
stress. It is therefore a strong nominal / single-goal variant or ablation,
while Risk Adapter v1 is the current balanced learned / risk-conditioned
protagonist.

\subsection{Baselines and Method Variants}

The evaluation compares the governor against no-op, fixed, heuristic, and
learned / risk-conditioned variants: Original, Fixed Scale 0.85, Wind-Level
0.85, Fixed Scale 0.80, Risk Adapter v1, Risk Adapter v2, and Risk Adapter
v2.1. Wind-Level 0.85 is the single-goal specialist heuristic. Fixed Scale 0.80
is the goal-reissue stress specialist static frontier. Risk Adapter v1 is the
balanced learned / risk-conditioned protagonist, while Risk Adapter v2.1 is a
strong nominal variant that is weaker under stress.

\subsection{Notation Block}

Let $u_{\mathrm{ref}}(t)$ denote the nominal reference or command produced by
the planner/MPC stack. The execution governor outputs a speed scale $s_v$ and
an acceleration scale $s_a$, with $s_v \in (0,1]$ and $s_a \in (0,1]$. The
executed command is modulated by these scale factors through the
command-adaptation interface:

\begin{equation}
u_{\mathrm{exec}}(t) = A(u_{\mathrm{ref}}(t), s_v(t), s_a(t)),
\end{equation}

where $A$ denotes the existing command-adaptation operation. The
risk-conditioned governor is written abstractly as:

\begin{equation}
g(r_3(t), r_5(t), c(t)) \rightarrow (s_v(t), s_a(t)),
\end{equation}

where $c(t)$ is execution context. This notation describes the interface only;
it does not imply formal safety, stability, or optimality guarantees.

\subsection{Claim Boundary}

The execution governor is an empirical runtime adaptation layer. It is not a
formal safety filter, does not provide a safety guarantee, and does not replace
the planner, payload MPC, or SO(3) controller. Its claim is limited to improving
balanced protocol-level behavior under the tested strong-wind evaluation.
